\documentclass[12pt]{article}\pagestyle{empty}
\begin{document}
\section*{Arbeitsblatt: Lineare Funktionen --- Klasse 8}
A linear function has the form $y = mx + b$. The number $m$ is called the
slope and the number $b$ is called the y-intercept.

\subsection*{1. The slope}
The slope $m$ tells you how steep the line is. It is computed as the rise
divided by the run: $m = \frac{\Delta y}{\Delta x}$.

\subsection*{2. The y-intercept}
The number $b$ is the value of $y$ where the line crosses the vertical axis,
that is, when $x = 0$.

\subsection*{3. From two points to an equation}
Given two points, first compute $m$ using the rise over run, then substitute
one point into $y = mx + b$ and solve for $b$.

\textbf{Common student error:} many students believe a steeper line always has
a larger value of $b$, confusing steepness with height.
\end{document}
